% -------------
% Módulo con los paquetes default de la plantilla.
% -------------

% Bibliografías y referencias
\usepackage[
backend=biber, % Recomendado, legacy : bibtex
style=numeric-comp, % numeric, apa, authoryear
sortlocale=author, % auto, nty (o variaciones de orden), none
]{biblatex}
    % \bibliographystyle{stylename} % apalike, ieeetr, plain, unsrt, alpha
\bibliography{src/proyecto/references.bib} % Archivo con las referencias
% Comillas para biblatex
\usepackage{csquotes}

% Uso para macros
\usepackage{xstring}

% Gráficas
\usepackage{graphicx}

% Headers y footers
\usepackage{fancyhdr}

% Código
\usepackage{listings}

% Colores (colores dvisps)
% \usepackage[dvipsnames]{xcolor}

% pdflatex highlighting
% \usepackage{soul}

% Cálculos JIT
% \usepackage{calc}
% Small bold/TT caps
% \usepackage{bold-extra}
% Justificación del texto
\usepackage{ragged2e}

% Símbolos y fuenets matemáticas
\usepackage{amsmath,amssymb,amsfonts,mathrsfs,commath}

% Unidades del SisInt (útil con floats)
% \usepackage{siunitx}

% Sistemas de ecuaciones
% \usepackage{systeme}

% Multicolumnas en tablas
\usepackage{multicol} 
% Multirenglones en tablas
\usepackage{multirow}

% Para tablas estiladas
\usepackage{nicematrix}
% Tablas multipágina
\usepackage{xltabular}
% Tablas simples pero más bonitas (sin verticales)
\usepackage{booktabs}

% Lorem Ipsum
\usepackage{lipsum}

% Texto pegado a imagenes
\usepackage{wrapfig}

% Evita usar // y separa automáticamente
\usepackage{parskip}

% Format enumerates/itemize
\usepackage{enumitem}

% Format captions
\usepackage{caption}

% Captions distintos en misma figura
\usepackage{subcaption}

% Referenicias (in document)
\usepackage{hyperref}

% Float (opción H "h pero obligado")
\usepackage{float}

% Interlineado
\usepackage{setspace}

% Formatear verbatim
\usepackage{fancyvrb}

% Personalizar headings
\usepackage{titlesec}
